\documentclass[12pt]{article}
\title{Lab 2 Note}
\author{Lab Tutor}
\date{March 2, 2026}

\usepackage{graphicx}



\begin{document}
\maketitle

\tableofcontents

\vskip 2cm 

\listoftables
\listoffigures

\newpage



\section{Table}
\subsection{tabular environment}

Example 1:\\
\begin{tabular}{|c|c|}
	\hline
	True & False \\
	\hline
	logical 1 & logical 0 \\
	\hline
\end{tabular}

\vspace{5mm}

\noindent Example 2:\\
\begin{center}
\begin{tabular}{c|c}
	$f(x)$ & $f'(x)$ \\
	\hline
	$x^{n}$ & $nx^{n-1}$ \\
	$\sin x$ & $\cos x$
\end{tabular}
\end{center}

\subsection{table environment}
Positioning a table is easy if they're inside a \underline{float table environment}. We can also add captions to tables.\\

\noindent Example 3:\\
\begin{table}[h]
	\centering
	\begin{tabular}{|l|l|}
		\hline
		True & False \\
		\hline
		logical 1 & logical 0 \\
		\hline
	\end{tabular}
	\caption{Simple Table}
\end{table}

\noindent Example 4:\\
\begin{table}[h]
	\centering
	\begin{tabular}{|c|c|c|c|}
		\hline
		$p$ & $q$ & $p \vee q$ & $p \wedge q$\\
		\hline
		1 & 1 & 1 & 1\\
		1 & 0 & 1 & 0\\
		\hline
	\end{tabular}
	\caption{Truth Table}
\end{table}


\section{Figure}
Extra package \texttt{graphicx} is required for inserting images to \LaTeX\,files. To declare the use of this package, include the following command in the preamble:
\begin{center}
	\verb|\usepackage{graphicx}|
\end{center}

\noindent Example 5:\\
Here is one image captured \footnote{Tips: Press \textbf{Windows Key + Shift + S} to capture the screen on Windows system.} from the website:  https://www.latex-project.org/
\includegraphics[scale=0.6]{latex}

Similar to the table environment, we can also use a float figure environment to edit the position/caption of the image.\\

\noindent Example 6:\\
Plot the graphs of two functions $\sin x$ and $\cos x$ from 0 to $2\pi$.
\begin{figure}[h]
	\centering
	\includegraphics[width=0.8\textwidth]{function image}
	\caption{Python Plot}
\end{figure}

\emph{\footnotesize Note: the image files must be put into the same folder where the .tex file is saved.}

\section{Text Font}
\subsection{font styles}

Common font styles used in \LaTeX:\\

\noindent Example 7:\\
normal text\\
\textbf{boldface font}\\
\textit{italics font}\\
\textsc{smallcaps font}\\
\texttt{typewriter font}\\
\textsf{serif font}\\
\underline{underline font}\\

\subsection{font sizes}

Different font sizes supported in \LaTeX:\\

\noindent Example 8:\\
Here is a {\tiny sample word}.\\
Here is a {\scriptsize sample word}.\\
Here is a {\footnotesize sample word}.\\
Here is a {\small sample word}.\\
Here is a sample word in normal size.\\
Here is a {\large sample word}.\\
Here is a {\Large sample word}.\\
Here is a {\LARGE sample word}.\\
Here is a {\huge sample word}.\\
Here is a {\Huge sample word}.\\





\end{document}